%! Carrying Out a Test for a Population Mean — Unit 7, Lesson 7.5 — Exit Ticket (Answer Key)
\documentclass[10pt]{article}
\usepackage{apstats-article}
\usepackage{apstats-key}   % pulls in -boxes; do NOT also load -boxes
\newcommand{\TallMath}[1]{$\displaystyle #1\rule[-1.4em]{0pt}{3.2em}$}

\begin{document}

\pageheader{Unit 7, Lesson 7.5}{Exit Ticket --- Answer Key}
\namedateperiod

\vspace{0.15in}

A physical education teacher claims the mean number of push-ups students can complete is
24. A fitness researcher believes students can do \emph{more} than 24 push-ups on average.
She randomly tests 18 students. $\bar x = 27.1$, $s = 6.3$. The dotplot is roughly symmetric
with no outliers. All conditions have been verified.

\begin{enumerate}
  \item Compute the $t$-statistic. Show your work.

        $SE = \dfrac{s}{\sqrt{n}} = $ \ans{$6.3/\sqrt{18} \approx 1.485$} \qquad
        $t = \dfrac{\bar x - \mu_0}{SE} = $ \ans{$(27.1 - 24)/1.485$} $\approx$ \ans{$2.09$}

  \item Find the $p$-value. Circle the correct TI-84 command, then give the result.

        \texttt{tcdf($-10^{99}$, $t$, df)} \quad OR \quad \ans{\texttt{tcdf($t$, $10^{99}$, df)}} \quad
        $p \approx$ \ans{$0.026$}

        \begin{teachernote}
        \texttt{tcdf(2.09, $10^{99}$, 17) $\approx$ 0.026}. Right-tail because $H_a: \mu > 24$.
        \end{teachernote}

  \item At $\alpha = 0.05$, write a complete conclusion in context.
        \ansline{Because $p \approx 0.026 < \alpha = 0.05$, we reject $H_0$. We have convincing evidence that the true mean number of push-ups students can complete is greater than 24.}
\end{enumerate}

\end{document}
